% Source: https://artofproblemsolving.com/community/c3h1671978_graphing_problem
% Stable identity: aops-topic-1671978-q01
Write the line as $2x+7y-8=0$. The shortest distance is along the perpendicular from $P(9,10)$ to the line. A normal vector of the line is $\mathbf n=(2,7)$, so the foot of the perpendicular is $F=P-t\mathbf n=(9-2t,\ 10-7t)$ for the $t$ with
\[
2(9-2t)+7(10-7t)=8\quad\Longleftrightarrow\quad88-53t=8\quad\Longleftrightarrow\quad t=\dfrac{80}{53}.
\]
Then $PF=t\,|\mathbf n|=\dfrac{80}{53}\sqrt{2^2+7^2}=\dfrac{80}{\sqrt{53}}$.

\begin{asy}
usepackage("amsmath");
size(6cm);
defaultpen(fontsize(10pt));
draw((-4,0)--(11,0),gray,Arrow(TeXHead));draw((0,-2)--(0,11),gray,Arrow(TeXHead));
label("$x$",(11,0),E,gray);label("$y$",(0,11),N,gray);
pair P=(9,10),Fp=(9-160/53,10-560/53);
draw((-4,16/7)--(11,-14/7),linewidth(0.9));
draw(P--Fp,red+dashed);
pair u=unit(P-Fp),v=unit((11,-2)-(-4,16/7));draw(Fp+0.5u--Fp+0.5u+0.5v--Fp+0.5v);
dot(P);dot(Fp);
label("$P(9,10)$",P,NE);label("$F$",Fp,SW);
label("$2x+7y=8$",(-4,16/7),NE);
label("$\dfrac{80}{\sqrt{53}}$",(P+Fp)/2,E,red);
label("\mbox{}",(min(currentpicture,user=true)+max(currentpicture,user=true))/2,black); // keeps the text after the figure black when LaTeX typesets the labels inline
\end{asy}

(This is the formula $\dfrac{|ax_0+by_0+c|}{\sqrt{a^2+b^2}}=\dfrac{|18+70-8|}{\sqrt{53}}$.) The distance is
\[
\boxed{\dfrac{80\sqrt{53}}{53}}\approx10.99 .
\]
Here $F=\left(\dfrac{317}{53},-\dfrac{30}{53}\right)$, and indeed $2\cdot317+7\cdot(-30)=424=8\cdot53$.
